\documentclass[10pt,a4paper]{amsart}
\usepackage{amsmath,amssymb,amsthm,graphicx}
% Retained minimal preamble for the Stacks Project chapter more-algebra.tex (More on Algebra).
% The chapter file pinned-original-source.tex carries only \input{preamble} on its first line, so no
% preamble text lives inside the chapter itself.  The Stacks Project uses one single shared preamble for
% every chapter; the snapshot stages it identically as sources/stacks-chow/preamble.tex and
% sources/stacks-etale-cohomology/preamble.tex.  The block below is copied mechanically, line for line,
% from sources/stacks-chow/preamble.tex lines 14-18 (the Xy-pic setup) and 154-217 (the native theorem
% declarations together with the mathematical macros).  Nothing was added, removed or redefined; the
% wrapper class and the amsmath/amssymb/amsthm/graphicx base packages are supplied by the assembler.
% Chapter: more-algebra.tex -- More on Algebra.
% For commutative diagrams we use Xy-pic
\usepackage[all]{xy}

% We use 2cell for 2-commutative diagrams.
\xyoption{2cell}

% Theorem environments.
%
\theoremstyle{plain}
\newtheorem{theorem}[subsection]{Theorem}
\newtheorem{proposition}[subsection]{Proposition}
\newtheorem{lemma}[subsection]{Lemma}

\theoremstyle{definition}
\newtheorem{definition}[subsection]{Definition}
\newtheorem{example}[subsection]{Example}
\newtheorem{exercise}[subsection]{Exercise}
\newtheorem{situation}[subsection]{Situation}

\theoremstyle{remark}
\newtheorem{remark}[subsection]{Remark}
\newtheorem{remarks}[subsection]{Remarks}

\numberwithin{equation}{subsection}

% Macros
%
\def\lim{\mathop{\mathrm{lim}}\nolimits}
\def\colim{\mathop{\mathrm{colim}}\nolimits}
\def\Spec{\mathop{\mathrm{Spec}}}
\def\Hom{\mathop{\mathrm{Hom}}\nolimits}
\def\Ext{\mathop{\mathrm{Ext}}\nolimits}
\def\SheafHom{\mathop{\mathcal{H}\!\mathit{om}}\nolimits}
\def\SheafExt{\mathop{\mathcal{E}\!\mathit{xt}}\nolimits}
\def\Sch{\mathit{Sch}}
\def\Mor{\mathop{\mathrm{Mor}}\nolimits}
\def\Ob{\mathop{\mathrm{Ob}}\nolimits}
\def\Sh{\mathop{\mathit{Sh}}\nolimits}
\def\NL{\mathop{N\!L}\nolimits}
\def\CH{\mathop{\mathrm{CH}}\nolimits}
\def\proetale{{pro\text{-}\acute{e}tale}}
\def\etale{{\acute{e}tale}}
\def\QCoh{\mathit{QCoh}}
\def\Ker{\mathop{\mathrm{Ker}}}
\def\Im{\mathop{\mathrm{Im}}}
\def\Coker{\mathop{\mathrm{Coker}}}
\def\Coim{\mathop{\mathrm{Coim}}}

% Boxtimes
%
\DeclareMathSymbol{\boxtimes}{\mathbin}{AMSa}{"02}

%
% Macros for moduli stacks/spaces
%
\def\QCohstack{\mathcal{QC}\!\mathit{oh}}
\def\Cohstack{\mathcal{C}\!\mathit{oh}}
\def\Spacesstack{\mathcal{S}\!\mathit{paces}}
\def\Quotfunctor{\mathrm{Quot}}
\def\Hilbfunctor{\mathrm{Hilb}}
\def\Curvesstack{\mathcal{C}\!\mathit{urves}}
\def\Polarizedstack{\mathcal{P}\!\mathit{olarized}}
\def\Complexesstack{\mathcal{C}\!\mathit{omplexes}}
% \Pic is the operator that assigns to X its picard group, usage \Pic(X)
% \Picardstack_{X/B} denotes the Picard stack of X over B
% \Picardfunctor_{X/B} denotes the Picard functor of X over B
\def\Pic{\mathop{\mathrm{Pic}}\nolimits}
\def\Picardstack{\mathcal{P}\!\mathit{ic}}
\def\Picardfunctor{\mathrm{Pic}}
\def\Deformationcategory{\mathcal{D}\!\mathit{ef}}
\usepackage{cleveref}
\crefname{originallemma}{original lemma}{original lemmas}
\Crefname{originallemma}{Original lemma}{Original lemmas}
\makeatletter
\newlabel{more-algebra-lemma-compare-completion-derived-completion}{{0H32}{source}{}{sourcepointer-1}{}}
\newlabel{more-algebra-lemma-compare-completion-derived-completion@cref}{{[originallemma][0][]0H32}{[1][0][]source}{}{}{}}
\makeatother
\crefname{originallemma}{original lemma}{original lemmas}
\Crefname{originallemma}{Original lemma}{Original lemmas}
\makeatletter
\newlabel{example-spectral-sequence-principal}{{0BKG}{source}{}{sourcepointer-2}{}}
\newlabel{example-spectral-sequence-principal@cref}{{[originallemma][0][]0BKG}{[1][0][]source}{}{}{}}
\makeatother
\crefname{originallemma}{original lemma}{original lemmas}
\Crefname{originallemma}{Original lemma}{Original lemmas}
\makeatletter
\newlabel{lemma-break-long-exact-sequence-modules}{{0CQE}{source}{}{sourcepointer-3}{}}
\newlabel{lemma-break-long-exact-sequence-modules@cref}{{[originallemma][0][]0CQE}{[1][0][]source}{}{}{}}
\makeatother
\crefname{originallemma}{original lemma}{original lemmas}
\Crefname{originallemma}{Original lemma}{Original lemmas}
\makeatletter
\newlabel{homology-lemma-Mittag-Leffler}{{02N1}{source}{}{sourcepointer-4}{}}
\newlabel{homology-lemma-Mittag-Leffler@cref}{{[originallemma][0][]02N1}{[1][0][]source}{}{}{}}
\makeatother
\title{The comparison diagram between adic completion and derived completion for a principal ideal}
\author{The Stacks Project Authors}
\date{stacks-project:more-algebra.tex:local-snapshot-20261012; GNU FDL 1.2 or later}
\begin{document}
\maketitle
\setcounter{subsection}{1}
\noindent\textit{Editorial note.} All delivered mathematical text is the author's own, copied from the pinned Stacks Project chapter file more-algebra.tex; no statement or proof step was rewritten, reordered or removed.\par\medskip
\noindent\textit{Editorial note.} The chapter file pinned-original-source.tex carries only \textbackslash{}input\{preamble\} on its first line, so no preamble text lives inside the chapter. The retained file sources/stacks-more-algebra/retained-preamble.tex is a mechanical, line-for-line copy of the project's single shared preamble as staged in the snapshot (Xy-pic setup plus theorem declarations and mathematical macros). Nothing was added, removed or redefined, and the excerpt serves only as preamble, never as a span coordinate.\par\medskip
\noindent\textit{Editorial note.} Both counted spans start at the \textbackslash{}begin line of their environment, so the statement and the author's proof are delivered whole; surrounding prose of the section is deliberately outside the counted spans and is carried as standard prerequisites instead.\par\medskip
\noindent\textit{Editorial note.} The citation fields (source\_id, source\_version, authors, source\_url, license.name, source\_title) are identical to those of every other card this lane has produced from this chapter, so the batch selector treats them as one source.\par\medskip
\noindent\textit{Allowed standard prerequisites (editorial).}\begin{itemize}
\item Derived completion of a complex along a principal ideal, and the complexes K\_n obtained by tensoring with the map f\textasciicircum{}n. used throughout the proof, lines 26959-27000. Supplies the two objects being compared.
\item Inverse limits of short exact sequences and the Mittag-Leffler condition. the author cites Homology, Lemma homology-lemma-Mittag-Leffler at proof lines 26970-26972. Justifies passing to the inverse limit and keeping exactness.
\item The long exact sequence associated with a principal ideal and the f-power torsion T\_f. used at proof lines 26980-27000. Supplies the third column of the diagram and the identification of the boundary terms.
\end{itemize}\medskip
\ifdefined\hypertarget\hypertarget{sourcepointer-1}{}\fi
\noindent\textit{Source reference.} Original lemma 0H32: focused-seeds/AGP/sources/stacks-snapshot-20261011/tags.txt, source lines 21302--21302. The lemma's own canonical tag, quoted so that readers can resolve the label inside the delivered text; it points at the counted statement itself and carries no proof content.\par\medskip
\ifdefined\hypertarget\hypertarget{sourcepointer-2}{}\fi
\noindent\textit{Source reference.} Original lemma 0BKG: focused-seeds/AGP/sources/stacks-snapshot-20261011/tags.txt, source lines 14559--14559. Navigation pointer to a statement the counted proof refers to but does not contain; nothing load-bearing is delegated to it.\par\medskip
\ifdefined\hypertarget\hypertarget{sourcepointer-3}{}\fi
\noindent\textit{Source reference.} Original lemma 0CQE: focused-seeds/AGP/sources/stacks-snapshot-20261011/tags.txt, source lines 15959--15959. Navigation pointer to a statement the counted proof refers to but does not contain; nothing load-bearing is delegated to it.\par\medskip
\ifdefined\hypertarget\hypertarget{sourcepointer-4}{}\fi
\noindent\textit{Source reference.} Original lemma 02N1: focused-seeds/AGP/sources/stacks-snapshot-20261011/tags.txt, source lines 3464--3464. Navigation pointer to a statement the counted proof refers to but does not contain; nothing load-bearing is delegated to it.\par\medskip
\begin{lemma}
\label{lemma-compare-completion-derived-completion}
Let $A$ be a ring and let $f \in A$. Let $K$ be an object of $D(A)$.
Denote $K_n = K \otimes_A^\mathbf{L} (A \xrightarrow{f^n} A)$.
For all $p \in \mathbf{Z}$ there is a commutative diagram
$$
\xymatrix{
& 0 & 0 \\
0 \ar[r] &
\widehat{H^p(K)} \ar[r] \ar[u] &
\lim H^p(K_n) \ar[r] \ar[u] &
T_f(H^{p + 1}(K)) \ar[r] &
0 \\
0 \ar[r] &
H^0(H^p(K)^\wedge) \ar[r] \ar[u] &
H^p(K^\wedge) \ar[r] \ar[u] &
T_f(H^{p + 1}(K)) \ar[r] \ar@{=}[u] &
0 \\
&
R^1\lim H^p(K)[f^n] \ar[u] \ar[r]^\cong &
R^1\lim H^{p - 1}(K_n) \ar[u] \\
& 0 \ar[u] & 0 \ar[u]
}
$$
with exact rows and columns where $\widehat{H^p(K)} = \lim H^p(K)/f^nH^p(K)$
is the usual $f$-adic completion. The left vertical short exact sequence
and the middle horizontal short exact sequence are taken from
Example \ref{example-spectral-sequence-principal}
The middle vertical short exact sequence is the one from
Lemma \ref{lemma-break-long-exact-sequence-modules}.
\end{lemma}

\begin{proof}
To construct the top horizontal short exact sequence, observe that
we have the following inverse system short exact sequences
$$
0 \to H^p(K)/f^nH^p(K) \to H^p(K_n) \to H^{p + 1}(K)[f^n] \to 0
$$
coming from the construction of $K_n$ as a shift of the
cone on $f^n : K \to K$. Taking the inverse limit of these
we obtain the top horizontal short exact sequence, see
Homology, Lemma \ref{homology-lemma-Mittag-Leffler}.

\medskip\noindent
Let us prove that we have a commutative diagram as in the lemma.
We consider the map $L = \tau_{\leq p}K \to K$.
Setting $L_n = L \otimes_A^\mathbf{L} (A \xrightarrow{f^n} A)$
we obtain a map $(L_n) \to (K_n)$ of inverse systems which
induces a map of short exact sequences
$$
\xymatrix{
0 & 0 \\
\lim H^p(L_n) \ar[r] \ar[u] &
\lim H^p(K_n) \ar[u] \\
H^p(L^\wedge) \ar[r] \ar[u] &
H^p(K^\wedge) \ar[u] \\
R^1\lim H^{p - 1}(L_n) \ar[r] \ar[u] &
R^1\lim H^{p - 1}(K_n) \ar[u] \\
0 \ar[u] & 0 \ar[u]
}
$$
Since $H^i(L) = 0$ for $i > p$ and $H^p(L) = H^p(K)$,
a computation using the references in the statement of the lemma
shows that $H^p(L^\wedge) = H^0(H^p(K)^\wedge)$
and that $H^p(L_n) = H^p(K)/f^nH^p(K)$. On the other hand,
we have $H^{p - 1}(L_n) = H^{p - 1}(K_n)$ and hence we see that
we get the isomorphism as indicated in the statement of the lemma
since we already know the kernel of
$H^0(H^p(K)^\wedge) \to \widehat{H^p(K)}$ is equal
to $R^1\lim H^p(K)[f^n]$. We omit the verification
that the rightmost square in the diagram commutes if we
define the top row by the construction in the first paragraph
of the proof.
\end{proof}

\end{document}
